\section{Localization, Boundaries, and Systems Proxy}
\label{sec:boundaries}

\subsection{Localization and candidate recovery}

E19 removes two oracle assumptions one at a time in a fixed-slot random-code
setting.  Separate-address control gains
$\CATENAResult{E19_SEPARATE_ADDRESS_GAIN}$ in the learned-address condition;
state-aware control gains $\CATENAResult{E19_STATE_READ_GAIN}$ when the erase
candidate must be read from current state; and the full controller gains
$\CATENAResult{E19_FULL_ONLY_GAIN}$ when both bottlenecks are active.  The
registered capable address/candidate floors and retention guards passed.  This
is learned localization over a fixed codebook, not semantic localization or
novel-entity generalization.

\subsection{Closed semantic and official claims}

E05a-R1 repaired relation representation, but the parameter-matched shared
controller also approached the oracle neighborhood.  Its registered absolute
factorized-advantage and direction gate did not pass.  H5 is therefore closed
for this submission---terminated, not refuted---and no E05b claim is made.

\subsection{Sequence localization transfer did not open}

E21 tested learned localization and current-state erase-candidate recovery
over repeated structured events with a fixed identifier schema and explicit
demand/provenance fields.  The original E21b aggregate is frozen as
\textsc{inconclusive-gate-implementation}: its guardrail implementation did
not identify the intended active state-read and cellwise non-inferiority
tests.  A prospectively locked aggregate-only repair retained the source runs
and primary estimands.

In E21b-R1, each of the three primary contrasts passed its registered effect
and 5/5 paired-seed direction gates.  The complete conjunction nevertheless
failed the capable affected-error floor and cellwise non-target
non-inferiority gate, so the frozen status is \textsc{not-supported}.  We
therefore treat the positive primary contrasts only as a negative boundary,
not as a selective-recovery or sequence-transfer claim.  E21 provides no
semantic or natural-language, novel-identifier, pretrained or recurrent-LM,
agent/planning, official-backend, or runtime-transfer evidence.

\subsection{Official and proxy boundaries}

E14 supports only a structured synthetic plan-state proxy.  E15a-R1 executed
pinned official GDN2 and FLA kernels, but tied-reduction and BF16 parity
exceeded their preregistered tolerances; the registered outcome remains
\textsc{fail}.  The no-patch audit attributed the tied divergence to the
official kernels' different IEEE-versus-TF32 solve policies at the first
inter-subchunk boundary, rather than probe wiring.  No reference fallback was
used.  E15b KVEraser remains unconfigured.  Consequently, no official
architecture-transfer, language-model, KVEraser/Transformer, or agent claim
is open.

\subsection{Quality-constrained systems proxy}

E20 compares one-time device-resident assimilation with three in-process
external-state proxies over a registered follow-up-query grid.  The first
quality-valid break-even query count is
$m=\CATENAResult{E20_BREAK_EVEN_EXTERNAL_PER_QUERY}$ for per-query external
state, $m=\CATENAResult{E20_BREAK_EVEN_COMPACT_CACHE}$ for the compact-cache
proxy, and $m=\CATENAResult{E20_BREAK_EVEN_FULL_REFRESH}$ for full refresh.
The internal path measures
$\CATENAResult{E20_INTERNAL_M1_MICROSECONDS}\,\mu$s at $m=1$ and
$\CATENAResult{E20_INTERNAL_M64_MICROSECONDS}\,\mu$s at $m=64$.

Persistent internal-state placement is outside the timed region; its
assimilation update is inside.  External baselines include their registered
CPU-to-device transfers.  This is a controlled in-process systems proxy, not a
storage-service or production-serving benchmark.

\subsection{Implication}

The controlled algebra addresses execution capacity once a demand is
specified.  Inferring that demand from semantics, localizing novel entities,
and transferring the mechanism to a quality-controlled official system remain
separate problems.
